% TODO(team): Keep setup and metric boundaries explicit during presentation.
\section{Method}
\begin{frame}{Independent processes on one shared host}
\small
\begin{tabularx}{\textwidth}{lY}
\toprule Controlled setting & Executed values \\\midrule
Topology & 3 Controllers + 3 Peers; real HTTP/TCP\\
Fixture & 9 distinct 128 KiB segments; 2 logical seconds each\\
Source limits & 96 / 144 / 192 KiB/s\\
Demand & Low: 2 viewers; high: 6; burst arrivals at 0 / 4 / 12 s\\
Client & Window 3; concurrency 2; startup buffer 4 s\\
Timeouts & 8 s per download attempt; 42 s per session\\\bottomrule
\end{tabularx}
\takeaway{Real byte transfer with simulated media consumption, not browser measurement.}
\note{Processes share CPU, RAM, disks and the host network stack. Fresh stores do not flush the OS page cache. Placement is fixed; source selection is round-robin.}
\end{frame}
\begin{frame}{Measure delivery and redundancy together}
\begin{columns}[T,onlytextwidth]
\column{.47\textwidth}
\begin{block}{Executed study}
81 unique runs in 67.9 minutes\\
72 E1--E4 runs + 9 E5 runs\\
396 scheduled viewer outcomes\\
Pilots excluded from estimates
\end{block}
\column{.49\textwidth}
\begin{block}{Outcomes}
Completion and observed stalls\\
Verified two-copy restoration\\
Repair payload\\
One-copy exposure (MiB$\cdot$s)
\end{block}
\end{columns}
\small First fault: 5 s. Observation horizon: 50 s. Recovery target: 40 s after fault. Three paired contexts; two per E4 cell.
\note{Treatment order is randomized and traces are paired. Failed viewers stay in the denominator. Inventory verification is sampled approximately every two seconds, so restoration is bracketed rather than exact.}
\end{frame}
